\section{Implementation}
\label{sec:implementation}

\subsection{Overview}

This section walks through the actual build of the lab: the network, the SIEM
manager, the two endpoints, the endpoint telemetry enhancements, and the
verification that logs flow end to end. It corresponds to the full
step-by-step documentation stored in the project repository under
\texttt{docs/}, distilled here into narrative form. The precise commands,
configuration snippets, and screenshots referenced in this section are held in
those source documents so this report can remain readable while still being
independently reproducible.

\subsection{Network and Virtualisation}
\label{sec:impl-network}

Oracle VirtualBox 7.2.8 was used to host all virtual machines. A dedicated
\textit{NAT network} named \texttt{SOC-Lab} was created on the address range
\texttt{10.0.10.0/24} as described in Section~\ref{sec:env-design}. The NAT
network mode was chosen deliberately over the more permissive Bridged mode:
VMs on a NAT network can reach each other but cannot be reached from, or
directly reach, the host network without deliberate port forwarding. This
provides the isolation guarantee that a security lab must offer by default.

Three virtual machines were provisioned onto this network:

\begin{itemize}[leftmargin=*]
    \item The Wazuh server VM at \texttt{10.0.10.3}.
    \item The Windows endpoint VM at \texttt{10.0.10.4}.
    \item The Linux endpoint VM at \texttt{10.0.10.5}.
\end{itemize}

Fixed addressing was preferred over DHCP for a lab of this size, so that
sources in the Wazuh dashboard would always resolve to the same host without
needing to cross-reference DHCP leases.

\subsection{Wazuh SIEM Deployment}
\label{sec:impl-wazuh}

The Wazuh server was deployed from the official
\texttt{wazuh-4.14.5.ova} appliance provided by Wazuh Inc. The appliance is
a pre-built all-in-one deployment containing the three Wazuh components
required to operate the platform:

\begin{itemize}[leftmargin=*]
    \item \textbf{wazuh-manager} --- the analysis and rule engine that receives
          agent telemetry, applies detection rules, and generates alerts.
    \item \textbf{wazuh-indexer} --- the storage and search layer, built on
          OpenSearch, that persists alerts and events so they can be queried
          and displayed.
    \item \textbf{wazuh-dashboard} --- the browser-based user interface, built
          on OpenSearch Dashboards, that presents alerts and enables analyst
          workflows.
\end{itemize}

Using the appliance rather than performing a from-source install was a
pragmatic choice: it reduced the initial deployment to VM import and network
attachment, which allowed the project to move on to the more instructive
work of endpoint configuration and detection engineering. The trade-off is
that the appliance is opaque compared to a from-source deployment; a
production build would follow Wazuh's documented distributed installation
procedure so component sizing could be tuned per role.

After boot, the appliance's default credentials were rotated, its network
interface was confirmed on \texttt{10.0.10.3}, and the dashboard was
reached at \texttt{https://10.0.10.3} from the host machine. The dashboard's
first-run wizard was used to generate agent deployment commands for the two
endpoints.

\subsection{Windows Endpoint}
\label{sec:impl-windows}

The Windows endpoint was provisioned from Microsoft's \texttt{WinDev2407Eval}
appliance (Windows 11), a free evaluation image intended for developer
testing. This choice avoided the licensing and installation overhead of a
manual Windows install while remaining fully representative of a modern
Windows endpoint.

Three configuration steps were performed after the VM was joined to the
\texttt{SOC-Lab} network:

First, the \textbf{Wazuh agent} was installed using the PowerShell one-liner
generated by the manager's Deploy new agent wizard. The wizard tailored the
command with the manager IP (\texttt{10.0.10.3}) and the assigned agent name
(\texttt{Windows-Endpoint}). The agent service was started with
\texttt{NET START WazuhSvc} and appeared as \texttt{Active} in the
manager's agent list within seconds.

Second, \textbf{Sysmon} was installed with the SwiftOnSecurity baseline
configuration. Sysmon writes to a dedicated Windows event channel,
\texttt{Microsoft-Windows-Sysmon/Operational}, which is separate from the
Windows Security event log and requires explicit inclusion in the Wazuh
agent's configuration to be forwarded. The following block was added to the
agent's \texttt{ossec.conf} inside the \texttt{<ossec\_config>} element:

\begin{quote}\small
\texttt{<localfile>}\\
\texttt{~~<location>Microsoft-Windows-Sysmon/Operational</location>}\\
\texttt{~~<log\_format>eventchannel</log\_format>}\\
\texttt{</localfile>}
\end{quote}

The agent service was restarted, and within one minute the manager began
receiving events under the \texttt{sysmon}, \texttt{sysmon\_event1}, and
\texttt{sysmon\_event\_3} rule groups. This confirmed that the SwiftOnSecurity
configuration was producing telemetry and that the Wazuh manager's Sysmon
decoder was correctly parsing the incoming events.

Third, a \textbf{baseline of normal activity} was observed for approximately
twenty-four hours. This served two purposes: it validated that the agent
would continue reporting stably across VM restarts, and it established a
reference of what "clean" telemetry from the endpoint looked like, so that
subsequent simulation-generated alerts could be distinguished from routine
noise.

Full step-by-step commands, screenshots, and verification outputs are
recorded in \texttt{docs/03-windows-endpoint.md} in the project repository.

\subsection{Linux Endpoint}
\label{sec:impl-linux}

The Linux endpoint was provisioned by performing a clean Ubuntu 26.04 LTS
Desktop install from the official ISO. Unlike Wazuh and Windows, both of which
were imported from vendor-provided appliances, the Linux install was carried
out manually. This was a deliberate choice: the Linux install exposed
first-hand the number of small decisions (partitioning, timezone, initial
user creation) that shape a system's behaviour and that an appliance import
hides.

The Wazuh agent was installed on the Ubuntu VM using the manager's
wizard-generated \texttt{wget} and \texttt{dpkg} command, then enabled and
started as a systemd unit. It appeared as \texttt{Active} in the dashboard
with \texttt{agent.id} 002.

The critical configuration step on Linux was adding \textbf{auditd
command-execution monitoring}. Wazuh's default Linux agent configuration
reads system logs from \texttt{/var/log/syslog}, \texttt{/var/log/auth.log},
and, if present, \texttt{/var/log/audit/audit.log}. However, Ubuntu Desktop
does not include the Linux Audit framework by default, so \texttt{auditd} was
installed via \texttt{apt}, and audit rules were added to track every
invocation of the \texttt{execve} system call --- which is called whenever
any command is executed on the system:

\begin{quote}\small
\texttt{-a exit,always -F arch=b64 -S execve -k audit-wazuh-c}\\
\texttt{-a exit,always -F arch=b32 -S execve -k audit-wazuh-c}
\end{quote}

The \texttt{-k audit-wazuh-c} key identifies these events as
command-execution audits, so they can be selected in Wazuh queries. After
the rules were loaded with \texttt{augenrules --load}, the Wazuh agent was
restarted and began forwarding \texttt{audit.log} entries alongside the
default log sources.

At this point, both endpoints were producing detection-grade telemetry, and
the milestone identified in the project overview as M4 (both endpoints
delivering enhanced telemetry to the SIEM) was met. Full commands, output,
and verification screenshots are recorded in
\texttt{docs/04-linux-endpoint.md}.

\subsection{Ingestion Verification and Baselining}
\label{sec:impl-verify}

Once both endpoints were reporting, three verification checks were performed
before any attack simulation was attempted. First, the number of events per
minute in the dashboard's \textit{Threat Hunting} view was observed for each
endpoint over a five-minute window and cross-referenced against known
activity on that endpoint (for example, opening a text editor on the Linux
endpoint produced a visible spike in \texttt{execve} events). Second, the
alert-level distribution was checked: a healthy Wazuh deployment shows a
long tail of Level~1 to Level~3 informational events, a small number of Level~5
to Level~8 alerts, and rare Level~12+ high-severity events, which was
observed. Third, a manual query for \texttt{rule.description: "Sysmon"} on
the Windows agent, and \texttt{rule.description: "execve"} on the Linux
agent, was used to confirm that both endpoints' enhanced telemetry was not
merely being ingested but was actually reachable through the dashboard's
query interface.

With these verifications passed, the lab was considered operational and
ready for the attack simulations described in Section~\ref{sec:evaluation}.
